\section{Introduction}

\todo{
\begin{itemize}
    \item Present the context of the problem and explain its practical relevance.
    \item Provide background information necessary to understand the problem domain.
    \item Briefly review related work or existing approaches to similar problems.
    \item Clearly define the problem being addressed.
    \item State the main objectives of the study.
    \item Highlight the expected contributions of the work.
\end{itemize}
}

Chronic Kidney Disease (CKD) represents a significant public health challenge
characterized by the progressive loss of renal function. Early identification of
individuals at risk is essential for timely intervention and improved clinical
outcomes. In recent years, machine learning approaches have emerged as promising
tools for supporting disease prediction tasks through the analysis of structured
clinical and laboratory data. Supervised classification methods, in particular,
have demonstrated strong potential for identifying patterns associated with
chronic diseases and supporting data-driven healthcare analysis
\cite{obermeyer2016,deo2015}.

The development of reliable predictive models for CKD involves multiple stages
of the machine learning workflow, including exploratory data analysis,
preprocessing, model selection, hyperparameter optimization, and model
interpretation. In clinical datasets, preprocessing techniques such as class
balancing and dimensionality reduction are frequently employed to address issues
including class imbalance, noisy attributes, and feature redundancy
\cite{he2009imbalanced,chawla2002smote}.
Additionally, robust evaluation strategies are necessary to ensure fair model
comparison and reproducibility of experimental results.

Prior work has demonstrated the applicability of machine learning to CKD
prediction using the same dataset adopted in this study. Islam et al.
\cite{inproceedings} evaluated Naive Bayes, Random Forest, Logistic Regression,
Decision Stump, and Linear Regression models for CKD risk factor prediction,
establishing a baseline for algorithm comparison on this data. Building on this
prior work, the present study extends the evaluation to a broader set of
algorithms, including Support Vector Machines (SVM) and K-Nearest Neighbors
(KNN), and incorporates class balancing and dimensionality reduction strategies
not explored in that work. Furthermore, explainability methods such as SHAP and
LIME have gained importance in healthcare-related applications by improving
transparency and supporting interpretation of model behavior.

In this context, this work investigates the impact of different preprocessing
and model optimization strategies on CKD prediction using the UCI Chronic Kidney
Disease dataset \cite{risk_factor_prediction_of_chronic_kidney_disease_857}. The
proposed workflow includes exploratory data analysis, evaluation of multiple
classification algorithms, application of balancing and dimensionality reduction
techniques, hyperparameter optimization using nested cross-validation, and
interpretability analysis of the final model. Experiments were conducted using
Stratified 5-Fold Cross-Validation to ensure consistent evaluation conditions.

The main objectives of this study are: (i) to compare the predictive performance
of different supervised learning algorithms for CKD prediction; (ii) to evaluate
the impact of class balancing and dimensionality reduction techniques on model
performance; (iii) to investigate hyperparameter optimization strategies for the
best-performing model; and (iv) to analyze model interpretability using feature
importance, permutation importance, SHAP, and LIME. Together, these objectives
aim to contribute a reproducible experimental framework for CKD classification
that pairs rigorous comparative evaluation of diverse supervised learning
algorithms with transparent post-hoc model interpretation.
